%======================================================================
% Region-transport covariance and torus translation covariance
%======================================================================

\begin{theorem}[Transport of the complement-side blocked weight]%
    \label{thm:peps_regionComplementWeight_transport}
    \lean{TNLean.PEPS.regionComplementWeight_transport}
    \leanok
    \uses{def:peps_tensor_transport, def:peps_regionComplementTwoBlock}
    Let $\phi:G\simeq G'$ be a graph isomorphism. The blocked weight of the
    transported tensor $A^\phi$ over the complement of the image region, read at
    the image complement boundary configuration and the image complement physical
    configuration, equals the blocked weight of $A$ over $V\setminus R$:
    \begin{align}
        T_{A^\phi,V'\setminus\phi(R)}^{(\phi\cdot\eta)^c}(\phi\cdot\tau)
        &=T_{A,V\setminus R}^{\eta^c}(\tau).
        \notag
    \end{align}
\end{theorem}
\begin{proof}\leanok
    The image of the complement of $R$ is the complement of the image of $R$, and
    after this identification the blocked weight transports along $\phi$ by the
    change of summation variables on virtual configurations.
\end{proof}

\begin{theorem}[Transport preserves agreement away from a bond]%
    \label{thm:peps_sameAwayFromBond_regionBoundaryConfigMap}
    \lean{TNLean.PEPS.sameAwayFromBond_regionBoundaryConfigMap}
    \leanok
    \uses{thm:peps_regionComplementWeight_transport}
    Let $\phi:G\simeq G'$ be a graph isomorphism and let $f$ be a crossing edge
    of $R$. Two crossing-edge labels of $R$ agree away from $f$ if and only if
    their images under the boundary reindex induced by $\phi$ agree away from the
    image of $f$.
\end{theorem}
\begin{proof}\leanok
    The boundary reindex is a bijection on crossing edges that carries $f$ to its
    image, so it carries the crossing edges other than $f$ bijectively to the
    crossing edges other than the image of $f$. Agreement on the former family
    is therefore equivalent to agreement on the latter.
\end{proof}

\begin{theorem}[Covariance of the region insertion coefficient]%
    \label{thm:peps_regionInsertedCoeff_transport}
    \lean{TNLean.PEPS.regionInsertedCoeff_transport}
    \leanok
    \uses{def:peps_regionInsertedCoeff, def:peps_tensor_transport,
        thm:peps_regionComplementWeight_transport,
        thm:peps_sameAwayFromBond_regionBoundaryConfigMap}
    Let $\phi:G\simeq G'$ be a graph isomorphism and let $f$ be a crossing edge
    of $R$. The region insertion coefficient of the transported tensor $A^\phi$
    over the image region, with $M$ inserted on the image of $f$ and the physical
    configurations relabelled by $\phi$, equals the region insertion coefficient
    of $A$ over $R$, with $M$ carried back to the bond of $A$ at $f$:
    \begin{align}
        C_{A^\phi}(\phi(R),\phi(f),M;\phi\cdot\sigma,\phi\cdot\tau)
        &=C_A(R,f,\phi^{*}M;\sigma,\tau).
        \notag
    \end{align}
    No endpoint orientation enters: the matrix is inserted between the two single
    boundary values on the edge.
\end{theorem}
\begin{proof}\leanok
    The two crossing-label sums reindex by the boundary action of $\phi$, the
    agreement-away-from-$f$ constraint is preserved by
    Theorem~\ref{thm:peps_sameAwayFromBond_regionBoundaryConfigMap}, the region
    weight transports along $\phi$, and the complement weight transports by
    Theorem~\ref{thm:peps_regionComplementWeight_transport}. The inserted matrix
    entry is read at the two boundary values, which the bond-dimension
    identification carries to the corresponding entry of $\phi^{*}M$.
\end{proof}

\begin{theorem}[Region insertion coefficient under an equality of tensors]%
    \label{thm:peps_regionInsertedCoeff_congrTensor}
    \lean{TNLean.PEPS.regionInsertedCoeff_congrTensor}
    \leanok
    \uses{def:peps_regionInsertedCoeff}
    If two tensors are equal, then their region insertion coefficients are equal,
    with the inserted matrix carried across the induced equality of bond
    dimensions:
    \begin{align}
        A_1=A_2
        \quad\Longrightarrow\quad
        C_{A_1}(R,f,M;\sigma,\tau)
        &=C_{A_2}(R,f,\psi M;\sigma,\tau).
        \notag
    \end{align}
    This applies to any tensor equality $A_1=A_2$, including the equality between
    a translation-invariant tensor and its transported copy.
\end{theorem}
\begin{proof}\leanok
    Substituting the equality of tensors identifies the two coefficients, and the
    induced equality of bond dimensions is the identity, so the carried matrix is
    the original one.
\end{proof}

\begin{theorem}[Bond dimension at a translated crossing edge]%
    \label{thm:peps_bondDim_boundaryEdgeMap_translate}
    \lean{TNLean.PEPS.bondDim_boundaryEdgeMap_translate}
    \leanok
    \uses{def:peps_torus_translation_invariant,
        def:peps_torus_geometry_translation}
    Let $A$ be translation invariant on the discrete torus, let $f$ be a crossing
    edge of $R$, and let $T_{a,b}$ be a translation. The bond dimension of $A$
    at the translated crossing edge equals its bond dimension at the reference
    crossing edge: $D_A(T_{a,b}(f))=D_A(f)$.
\end{theorem}
\begin{proof}\leanok
    Translation carries the reference edge to its image, and translation
    invariance fixes the tensor under transport, so the bond dimension is
    unchanged.
\end{proof}

\begin{theorem}[Translation covariance of the coefficient identity]%
    \label{thm:peps_regionInsertedCoeff_translate_coeffIdentity}
    \lean{TNLean.PEPS.regionInsertedCoeff_translate_coeffIdentity}
    \leanok
    \uses{thm:peps_regionInsertedCoeff_transport,
        thm:peps_regionInsertedCoeff_congrTensor,
        thm:peps_bondDim_boundaryEdgeMap_translate,
        def:peps_torus_translation_invariant}
    Let $A$ and $B$ be translation invariant on the discrete torus, and let $f$ be
    a crossing edge of $R$. Suppose a map $g$ realizes the coefficient identity
    at the reference edge:
    $C_A(R,f,M;\sigma,\tau)=C_B(R,f,g(M);\sigma,\tau)$.
    Then, for every translation $T_{a,b}$, the same identity holds at the
    translated crossing edge of the translated region, with the inserted matrix
    carried between the translated and reference bonds:
    \begin{align}
        C_A\bigl(T_{a,b}(R),T_{a,b}(f),M';
            T_{a,b}\cdot\sigma,T_{a,b}\cdot\tau\bigr)
        &=
        C_B\bigl(T_{a,b}(R),T_{a,b}(f),
            \widehat{g}(M');
            T_{a,b}\cdot\sigma,T_{a,b}\cdot\tau\bigr),
        \label{eq:peps_translate_coeff}
    \end{align}
    where $\widehat{g}$ applies $g$ to the reference-bond reindex of $M'$ and
    carries the result back to the translated bond.
\end{theorem}
\begin{proof}\leanok
    Translation invariance identifies the transported tensor with the original,
    so Theorem~\ref{thm:peps_regionInsertedCoeff_transport} carries the reference
    coefficient of $A$ to the translated edge with no change of tensor;
    Theorem~\ref{thm:peps_regionInsertedCoeff_congrTensor} mediates the literal
    and transported tensors. Applying the reference identity for $g$ and
    transporting back for $B$ in the same way gives
    (\ref{eq:peps_translate_coeff}), the inserted matrix being carried across
    the equal bond dimensions of
    Theorem~\ref{thm:peps_bondDim_boundaryEdgeMap_translate}.
\end{proof}

\begin{definition}[Single-vertex two-block tensor]%
    \label{def:peps_vertexTwoBlock}
    \lean{TNLean.PEPS.vertexTwoBlock}
    \leanok
    \uses{def:peps_tensor}
    Fix a vertex $v$. The single-vertex tensor $A_v$ is regarded as a
    two-block tensor whose shared bonds are the edges incident to $v$, whose
    external boundary is a one-point space, and whose physical index is the
    physical index at $v$: $T_v(*,\eta,\sigma)=A_v(\eta,\sigma)$.
\end{definition}

\begin{theorem}[Injectivity of the single-vertex two-block tensor]%
    \label{thm:peps_isTwoBlockInjective_vertexTwoBlock}
    \lean{TNLean.PEPS.isTwoBlockInjective_vertexTwoBlock}
    \leanok
    \uses{def:IsVertexInjective, def:peps_vertexTwoBlock}
    If $A$ is vertex-injective, then the single-vertex two-block tensor at every
    vertex is injective.
\end{theorem}
\begin{proof}\leanok
    Vertex injectivity is linear independence of the coefficient family
    $\eta\mapsto A_v(\eta,-)$. The auxiliary one-point boundary only reindexes
    this same family, so the two-block injectivity condition is the same linear
    independence statement.
\end{proof}

\begin{definition}[Complement two-block tensor]%
    \label{def:peps_complementTwoBlock}
    \lean{TNLean.PEPS.complementTwoBlock}
    \leanok
    \uses{def:peps_tensor}
    For a vertex $v$, the complementary block $V\setminus\{v\}$ defines a
    two-block tensor on the same incident shared bonds. Its coefficient is the
    contraction of all tensors away from $v$, with the incident bond indices
    left open at the complement endpoints:
    $T_{v^c}(*,\eta,\tau)=A_{v^c}(\eta,\tau)$.
\end{definition}

\begin{definition}[Complement-region index identifications]%
    \label{def:peps_vertexComplementRegionEquivs}
    \lean{TNLean.PEPS.isRegionBoundaryEdge_vertexComplementVertices_iff}
    \lean{TNLean.PEPS.vertexComplementBoundaryEdgeEquiv}
    \lean{TNLean.PEPS.vertexComplementBoundaryConfigEquiv}
    \lean{TNLean.PEPS.vertexComplementVertexEquiv}
    \lean{TNLean.PEPS.vertexComplementPhysicalConfigEquiv}
    \leanok
    \uses{def:peps_complementTwoBlock}
    For a vertex $v$, the edges crossing the boundary of
    $V\setminus\{v\}$ are precisely the edges incident to $v$. Hence the
    boundary configurations of $V\setminus\{v\}$ identify with the local
    virtual configurations at $v$. Similarly, the vertices in
    $V\setminus\{v\}$ identify with the vertices distinct from $v$, giving
    the corresponding identification of physical configurations.
\end{definition}

\begin{theorem}[Vertex-complement contraction as a blocked region]%
    \label{thm:peps_regionBlockedWeight_vertexComplement}
    \lean{TNLean.PEPS.vertexComplementBoundaryConfigEquiv_regionBoundaryLabel}
    \lean{TNLean.PEPS.regionBlockedWeight_vertexComplement}
    \leanok
    \uses{def:peps_vertexComplementRegionEquivs, def:peps_complementTwoBlock}
    Under the preceding identifications, the blocked tensor of
    $V\setminus\{v\}$ is the vertex-complement tensor:
    \begin{align}
        T_{A,V\setminus\{v\}}^{\rho}(\tau)
        &=
        A_{v^c}(\Psi_v\rho,\Phi_v\tau).
        \notag
    \end{align}
\end{theorem}
\begin{proof}\leanok
    Expanding the blocked contraction as a sum over the virtual
    configurations of the complement,
    \begin{align}
        T_{A,V\setminus\{v\}}^{\rho}(\tau)
        &=
        \textstyle
        \sum_{\zeta:\,\zeta|_{\partial}=\rho}\;
        \prod_{w\neq v} A_w(\zeta|_w,\tau_w)
        =
        A_{v^c}(\Psi_v\rho,\Phi_v\tau),
        \notag
    \end{align}
    where $\zeta$ runs over the virtual configurations of the complement
    whose boundary restriction $\zeta|_{\partial}$ is the prescribed boundary
    configuration $\rho$. The boundary identification $\Psi_v$ turns the
    crossing-edge constraint $\zeta|_{\partial}=\rho$ into the prescribed
    incident-edge labels, and the vertex identification $\Phi_v$ reindexes
    the product over the vertices distinct from $v$; both bijections match
    the summation sets and their factors termwise.
\end{proof}

\begin{theorem}[Injectivity of the vertex-complement contraction]%
    \label{thm:peps_vertexComplementTensorInjective}
    \lean{TNLean.PEPS.vertexComplementTensorInjective_of_isVertexInjective}
    \leanok
    \uses{def:IsVertexInjective, def:peps_tensor}
    If $A$ is vertex-injective and every virtual bond space has positive
    dimension, then the coefficient family
    $\eta\longmapsto(\tau\longmapsto A_{v^c}(\eta,\tau))$
    of the vertex-complement contraction is linearly independent.
\end{theorem}
\begin{proof}\leanok
    \uses{thm:peps_regionBlockedWeight_vertexComplement,
        thm:peps_regionBlockedTensorInjective_of_isVertexInjective}
    Apply the injectivity theorem for an arbitrary finite region to
    $R=V\setminus\{v\}$. Theorem~\ref{thm:peps_regionBlockedWeight_vertexComplement}
    identifies the resulting blocked tensor with the vertex-complement tensor.
    Reindexing the boundary family and the physical coordinates by bijections
    preserves linear independence.
\end{proof}

\begin{theorem}[Injectivity of the complement two-block tensor]%
    \label{thm:peps_isTwoBlockInjective_complementTwoBlock}
    \lean{TNLean.PEPS.isTwoBlockInjective_complementTwoBlock}
    \leanok
    \uses{def:IsVertexInjective, def:peps_complementTwoBlock,
        thm:peps_vertexComplementTensorInjective}
    If $A$ is vertex-injective and every virtual bond space has positive
    dimension, then the complement two-block tensor at every vertex is
    injective.
\end{theorem}
\begin{proof}\leanok
    \uses{thm:peps_vertexComplementTensorInjective}
    The previous theorem gives linear independence of the coefficient family
    $\eta\mapsto A_{v^c}(\eta,-)$. The two-block tensor only adds a one-point
    external boundary, so this is the same injectivity condition after the
    harmless reindexing $(*,\eta)\leftrightarrow\eta$.
\end{proof}

\begin{theorem}[One-vertex versus complement comparison]%
    \label{thm:peps_oneVertexComplementComparison}
    \lean{TNLean.PEPS.one_vertex_complement_comparison}
    \leanok
    \uses{thm:peps_postAbsorptionEdgeInsertionEquality,
        thm:peps_twoInjectiveTensorInsertionComparison}
    Let the shared bond type and the two external boundary spaces be nonempty.
    After the edge gauges have been absorbed into the second tensor family
    $\widetilde B$, block one vertex $v$ against its complement. Assume that
    $A_v,A_{v^c},\widetilde B_v,\widetilde B_{v^c}$ are injective as two-block
    tensors. The post-absorption insertion equality gives, for every shared
    bond $b$, every matrix $X\in\MN{D_b}$, and all external and
    physical indices,
    \begin{align}
        C_{A_v,A_{v^c}}(b,X;\eta_v,\eta_{v^c},\sigma_v,\sigma_{v^c})
        &=
        C_{\widetilde B_v,\widetilde B_{v^c}}
        (b,X;\eta_v,\eta_{v^c},\sigma_v,\sigma_{v^c}).
        \notag
    \end{align}
    Then $\exists\lambda_v\in\C^\times$ such that
    $A_v=\lambda_v\widetilde B_v$.
\end{theorem}
\begin{proof}\leanok
    With $A_1=A_v$, $A_2=A_{v^c}$, $B_1=\widetilde B_v$, and
    $B_2=\widetilde B_{v^c}$, the post-absorption equality is precisely, for
    every $b$, $X$, $\eta_v$, $\eta_{v^c}$, $\sigma_v$, and $\sigma_{v^c}$,
    \begin{align}
        C_{A_1,A_2}(b,X;\eta_v,\eta_{v^c},\sigma_v,\sigma_{v^c})
        &=
        C_{B_1,B_2}(b,X;\eta_v,\eta_{v^c},\sigma_v,\sigma_{v^c}).
        \notag
    \end{align}
    Theorem~\ref{thm:peps_twoInjectiveTensorInsertionComparison} gives
    \begin{align}
        \exists\lambda_v\in\C^\times,\qquad
        A_v&=\lambda_v\widetilde B_v,\qquad
        A_{v^c}=\lambda_v^{-1}\widetilde B_{v^c}.
        \notag
    \end{align}
    The desired local proportionality is the first equality.
\end{proof}

\begin{remark}[Normal-region form]
    The normal-region application in \cite[Section~4]{Molnar2018NormalPEPS}
    uses the same scalar comparison after comparing two injective regions
    $R\subset S=R\cup\{v\}$. If the region comparisons give
    $A_R=\lambda_R\widetilde B_R$ and
    $A_S=\lambda_S\widetilde B_S$, the source-paper comparison is
    % Formula: A_R A_v=A_S\propto\widetilde B_S=\widetilde B_R\widetilde B_v.
    % Source verdict: Molnar et al., arXiv:1804.04964, page 16, lines
    %   1544--1570 of paper_normal.tex.  The source labels the tensors at v
    %   simply by A and \widetilde B.
    % Ink-to-index map: the horizontal bond in each two-tensor panel is the
    %   contracted boundary index between R and v.  The upper stubs are free
    %   physical indices, and the outer horizontal stubs are free virtual
    %   boundary indices.
    % Type verdict: the single upper leg of A_S or \widetilde B_S denotes the
    %   fused physical space H_R\otimes H_v; the source draws that composite
    %   leg once, while the factored panels expose its H_R and H_v factors.
    % Contraction audit: the two-tensor panels have one internal virtual
    %   contraction; the one-tensor panels have none.
    % Two-tensor boundary signature:
    %   boundary|virtual-west=1|virtual-east=1|physical-up=2|physical-down=0.
    % One-tensor boundary signature:
    %   boundary|virtual-west=1|virtual-east=1|physical-up=1|physical-down=0.
    \begin{center}
        \tnpic[physical=up]{
            \tn[label pos=south]{A_R} & \tn[label pos=south]{A}
        }
        \quad$=$\quad
        \tnpic[physical=up]{
            \tn[label pos=south]{A_S}
        }
        \quad$\propto$\quad
        \tnpic[physical=up]{
            \tn[label pos=south]{\widetilde B_S}
        }
        \quad$=$\quad
        \tnpic[physical=up]{
            \tn[label pos=south]{\widetilde B_R} &
            \tn[label pos=south]{\widetilde B}
        }.
    \end{center}
    Applying the left inverse supplied by injectivity of the block
    $A_R=\lambda_R\widetilde B_R$ gives
    $A_v=(\lambda_S/\lambda_R)\widetilde B_v$.
\end{remark}

\begin{theorem}[Fundamental Theorem for injective PEPS, conditional on
        bond-dimension equality]%
    \label{thm:fundamentalTheorem_PEPS_of_bondDim}
    \lean{TNLean.PEPS.fundamentalTheorem_PEPS_of_bondDim}
    \leanok
    \uses{thm:gaugeConsistency, thm:peps_oneVertexComplementComparison,
        def:GaugeEquivPEPS, def:IsVertexInjective}
    Suppose that the corresponding virtual edge spaces of $A$ and $B$ are
    already identified, every virtual bond of $A$ has positive dimension, and
    the graph is connected.
    If $A$ and $B$ are vertex-injective and give the same PEPS state, then
    there are invertible matrices $X_e$, one for each edge, such that at every
    vertex the tensor $B_v$ is obtained from $A_v$ by the oriented endpoint
    action of the incident matrices.
\end{theorem}
\begin{proof}\leanok
    Gauge consistency supplies the edge gauges directly, and the gauge
    equivalence packages them with the identified bond dimensions.
\end{proof}

\begin{theorem}[Fundamental Theorem for injective PEPS]%
    \label{thm:fundamentalTheorem_PEPS}
    \lean{TNLean.PEPS.fundamentalTheorem_PEPS}
    \leanok
    % Scope restriction: this is the positive-bond variant of Theorem 2 in
    % \cite{Molnar2018NormalPEPS}. The elimination plan for the positivity
    % hypotheses is recorded in
    % \path{docs/paper-gaps/peps_injective_ft_section3_route.tex}.
    \uses{thm:fundamentalTheorem_PEPS_of_bondDim, def:GaugeEquivPEPS,
        def:IsVertexInjective}
    Let $A$ and $B$ be vertex-injective PEPS tensors on a finite simple
    graph, and suppose every virtual bond of $A$ and $B$ has positive
    dimension and that the graph is connected. If they give the same PEPS
    state, then their corresponding
    virtual edge spaces have the same dimensions and, after identifying these
    spaces, there are invertible matrices $X_e$, one for each edge, such that
    $B_v$ is obtained from $A_v$ by the oriented endpoint action of the matrices
    incident to $v$.
\end{theorem}
\begin{proof}\leanok
    The bond dimensions agree edgewise, because blocking around an edge gives
    two injective three-site chains whose matched matrix insertions on that bond
    form an algebra equivalence between the two bond matrix algebras, forcing
    equal sizes. With the bond dimensions identified, the conditional form of
    the theorem applies.
\end{proof}
